% No preamble. Prefix pgl-. Entire pinned ghost coefficient; vector still needed.
\section{The pinning scale and the entire based-ghost logarithm}
\label{pgl-section}

The pinning scale is longer than the block diffusion
time. We now control it sharply enough to determine
the whole actual based-ghost Maxwell contribution.
All statements retain the infinite-system and formal
jet conventions of Section~\ref{pgp-section}.

\subsection{The correct flat pinned gap}
\label{pgl-gap-section}

\begin{lemma}[Four-dimensional point-pinned Poincare bound]
\label{pgl-gap}
Let \(G_K(0)\) be the flat scalar lattice Laplacian on
functions vanishing at every site of \(K\mathbb Z^4\).
There is a constant \(c>0\), independent of \(K\), such
that
\begin{equation}
                     G_K(0)\ge cK^{-4}I.
 \label{pgl-gap-bound}
\end{equation}
The statement holds on compatible finite periodic
volumes and on the infinite lattice.
\end{lemma}

\begin{proof}
It suffices to prove a block inequality using only
edges internal to each centered cube. For \(K=1\)
the domain is zero; assume \(K\ge2\).
The Neumann graph Laplacian of a path with \(K\)
vertices has eigenvalues
\[
 \lambda_j=4\sin^2\frac{\pi j}{2K},\qquad
                  0\le j<K.
\]
Its normalized cosine eigenvectors have squared
values at most \(2/K\). Products give an eigenbasis
\(\psi_{\boldsymbol j}\) in four dimensions, with
\[
 |\psi_{\boldsymbol j}(x)|^2\le16K^{-4},\qquad
 \lambda_{\boldsymbol j}
       =\sum_{\ell=1}^4\lambda_{j_\ell}
                          \ge c|\boldsymbol j|^2/K^2.
\]
Consequently the point diagonal of the zero-mean
Neumann inverse is uniformly bounded:
\[
 \sum_{\boldsymbol j\ne0}
       \frac{|\psi_{\boldsymbol j}(x)|^2}
            {\lambda_{\boldsymbol j}}
 \le CK^{-2}
       \sum_{\substack{\boldsymbol j\in\{0,\ldots,K-1\}^4\\
                        \boldsymbol j\ne0}}
                              |\boldsymbol j|^{-2}\le C.
\]
The last sum is \(O(K^2)\), by shells in four dimensions.

Write a block field as \(\phi=\bar\phi+u\), with
\(u\) of mean zero. Its internal edge energy is
denoted by \(E_B(\phi)\). Spectral Cauchy--Schwarz
gives \(|u(x)|^2\le C E_B(\phi)\), and the first
nonzero Neumann eigenvalue gives
\(\|u\|_2^2\le CK^2 E_B(\phi)\).
At the pinned center, \(\bar\phi=-u(x_{\rm center})\).
It follows that
\[
 \|\phi\|_{2,B}^2=\|u\|_{2,B}^2+K^4|\bar\phi|^2
                                  \le CK^4 E_B(\phi).
\]
Sum over disjoint blocks. Internal edges are a
subset of all edges, proving \eqref{pgl-gap-bound}.
The infinite statement follows first on finitely
supported fields and then by closure.
\end{proof}

Together with the constant-off-centers trial function
from Notebook 32, this gives the sharp order \(K^{-4}\)
of the bottom of the flat based spectrum. The finite
periodic upper test gives at most \(8/(K^4-1)\);
the same order follows in infinite volume by large
cutoffs. It does not give a mass of order \(K^{-1}\).

\subsection{A killed-return bound and an area fourth moment}
\label{pgl-tail-inputs-section}

Let \(p_s^K(x,y)\) be the heat kernel killed on visiting
a center. At the flat background it is the kernel of
\(e^{-sG_K(0)}\). For \(s\ge1\), free-kernel domination
and \eqref{pgl-gap-bound} imply
\begin{equation}
 p_s^K(x,x)\le Cs^{-2}e^{-c s/K^4},
                         \qquad x\notin K\mathbb Z^4.
 \label{pgl-killed-return}
\end{equation}
Indeed split \(e^{-sG_K/2}\delta_x\) into a time
\(s/4\) contraction followed by time \(s/4\)
smoothing:
\[
 p_s^K(x,x)
 =\|e^{-sG_K/2}\delta_x\|_2^2
 \le e^{-c s/(2K^4)}p_{s/2}^K(x,x)
 \le e^{-c s/(2K^4)}p_{s/2}^{\rm free}(0,0).
\]
The last free return is \(O(s^{-2})\) by the
one-dimensional integral estimates already proved.
Absorb the harmless factor two in \(c\).

\begin{lemma}[Fourth area moment of the returning walk]
\label{pgl-area-fourth}
For \(s\ge1\), the free four-dimensional walk satisfies
\begin{equation}
 \mathbb E[{\cal A}(s)^4\mid X(s)=0]\le Cs^4,\qquad
 \mathbb E[{\cal A}(s)^4{\bf1}_{\{X(s)=0\}}]\le Cs^2.
 \label{pgl-fourth-bound}
\end{equation}
\end{lemma}

\begin{proof}
For a returning one-dimensional coordinate, its
balanced jump count \(m\) satisfies the exact identity
\[
 \mathbb E[m^2\mid X_1(s)=0]
 =\frac{\sum_{m\ge1}m^2s^{2m}/(m!)^2}{I_0(2s)}
 =s^2.
\]
Thus \(\mathbb E m\le s\) as well. Given the bridge
path \(Y=X_\nu\) and the count \(m\) for \(X_\mu\),
its area is
\[
 {\cal A}(s)=\sum_{i=1}^m Y(U_i)-\sum_{i=1}^m Y(V_i),
\]
where the \(U_i,V_i\) are independent uniforms in
\([0,s]\), independent of \(Y\). The deterministic
mean terms cancel. The fourth moment formula for a
sum of independent centered variables and Jensen's
inequality give
\[
 \mathbb E[{\cal A}(s)^4\mid m,Y]
 \le C(m+m^2)\,\frac1s\int_0^s|Y(u)|^4\,du.
\]
For clarity, centering \(Y(U_i)\) changes its fourth
moment by at most a fixed factor; the square of its
second moment is at most its fourth moment.

Given the balanced count \(l\) of \(Y\), the value
\(Y(u)\) is the difference of two independent
\(\operatorname{Binomial}(l,u/s)\) variables.
The same centered-sum formula bounds its fourth
moment by \(C(l+l^2)\), uniformly in \(u\).
Averaging counts therefore gives
\[
 \mathbb E[|Y(u)|^4\mid Y(s)=0]\le C(s+s^2)\le Cs^2.
\]
The two coordinate bridges and their counts are
independent. Combining the last bounds proves
the conditional \(Cs^4\) estimate. The four-dimensional
return probability is \(p_s^4\le Cs^{-2}\), proving
the unconditional bound.
\end{proof}

\begin{proposition}[Uniform ghost tail past the pinning scale]
\label{pgl-tail}
For \(s\ge1\),
\begin{equation}
 0\le h_K^g(s)\le C e^{-c s/K^4}.
 \label{pgl-ghost-tail-pointwise}
\end{equation}
In particular,
\begin{equation}
           \int_{K^4}^\infty h_K^g(s)\,\frac{ds}{s}\le C,
 \label{pgl-tail-integral}
\end{equation}
uniformly in \(K,M\).
\end{proposition}

\begin{proof}
At each noncenter starting site \(x\), the squared-area
walk formula has the probabilistic form
\[
 h_K^g(s)=\frac{C_H}{K^4}
 \sum_{\substack{x\ {\rm noncenter}\\x\ {\rm in\ block}\ 0}}
 \mathbb E_x\left[{\cal A}(s)^2
    {\bf1}_{\{X(s)=x,\ X\ {\rm avoids\ centers}\}}\right].
\]
Translate the area coordinate to \(x\), which changes
no area on a closed walk. Cauchy--Schwarz bounds each
expectation by
\[
 \left(\mathbb E[{\cal A}(s)^4
                       {\bf1}_{\{X(s)=0\}}]\right)^{1/2}
                  \left(p_s^K(x,x)\right)^{1/2}
 \le C e^{-c s/K^4},
\]
using \eqref{pgl-killed-return} and
\eqref{pgl-fourth-bound}. The block has at most \(K^4\)
sites, canceling the prefactor. Integrate with
\(s=K^4v\), \(v\ge1\), to obtain
\eqref{pgl-tail-integral}.
\end{proof}

\subsection{The additional logarithm that the vector must cancel}
\label{pgl-full-section}

\begin{theorem}[Entire actual based-ghost coefficient]
\label{pgl-full}
There is a constant \(C\), independent of \(M,K\) and
the permitted joint structural system, such that
\begin{equation}
 \left|\int_0^\infty h_K^g(s)\,\frac{ds}{s}
                      -4\alpha_g\log K\right|\le C.
 \label{pgl-full-law}
\end{equation}
More specifically,
\begin{equation}
 \int_{K^2}^{K^4} h_K^g(s)\,\frac{ds}{s}
                         =2\alpha_g\log K+O(1),
 \label{pgl-extra-log}
\end{equation}
and the tail past \(K^4\) is uniformly bounded.
\end{theorem}

\begin{proof}
The proof of Proposition~\ref{pgp-pinning} gave, for
every \(T>0\),
\[
 0\le\int_0^T(h_\infty^g-h_K^g)\,\frac{ds}{s}
 \le K^{-4}\left\{h_\infty^g(T)
                    +8\int_0^T h_\infty^g(s)\,ds\right\}.
\]
Set \(T=K^4\). The right side is uniformly bounded,
while the unpinned integral is
\(\alpha_g\log(K^4)+\gamma_g+O(K^{-4})\).
This proves the full prefix law with an \(O(1)\)
error. Add Proposition~\ref{pgl-tail} to get
\eqref{pgl-full-law}. Subtracting
\eqref{pgp-prefix-law} gives \eqref{pgl-extra-log}.
\end{proof}

Recall explicitly that
\[
 \int_0^\infty h_K^g(s)\,\frac{ds}{s}
   =2[t^2]\operatorname{Tr}_{0,S}\log G(t)
   =\beta\bigl(\mathcal N_0\log\det G\bigr),
\]
with the determinant-density and Maxwell normalization
of Notebook 33. Thus the displayed factor two is the
affine coefficient extraction, not an extra determinant.
Its coefficient \(4\alpha_g=C_H/(12\pi^2)\) is
twice the \(2\alpha_g=C_H/(24\pi^2)\) of the
continuum ghost with mass \(K^{-1}\) in current
fine units and the same ultraviolet scale.
This difference is an actual pinning effect in the
separated determinant; it is not a failure of the
normalization.

In particular, a uniform estimate on the separated
ghost tail starting at \(K^2\) is false. Its positivity
and \eqref{pgl-extra-log} prove a logarithmic lower
bound on that tail. The total reference coefficient
is \(22\alpha_g\log K\); hence a bounded comparison
for the full actual determinant requires the vector
contribution to be \(18\alpha_g\log K+O(1)\), since
the coarse determinant is bounded. Equivalently,
the vector must compensate the additional ghost
logarithm relative to its massive reference.
No such vector estimate is established by this
scalar calculation.
